\documentclass[11pt]{article}
\usepackage[a4paper,margin=27mm]{geometry}
\usepackage{amsmath,amssymb,amsthm,xcolor,booktabs,array}
\newcommand{\PROVED}{\textcolor{green!40!black}{\textbf{PROVED}}}
\newcommand{\OPEN}{\textcolor{red!70!black}{\textbf{OPEN}}}
\newcommand{\AUDIT}{\textcolor{orange!70!black}{\textbf{AUDIT}}}
\newtheorem{theorem}{Theorem}
\newtheorem{corollary}{Corollary}
\title{Uniform Gamma--Tate coupling at an arithmetic threshold cell (v112)}
\date{13 August 2026}

\begin{document}
\maketitle

\section{The off-diagonal Gamma kernel}

At the critical parameter write $a_j=2j+\tfrac12$.  The positive Gamma
energy has multiplier
\begin{equation}
 g_\Gamma(\tau)=\sum_{j\ge0}
 \frac{2\tau^2}{a_j(a_j^2+\tau^2)}.                            \tag{1}
\end{equation}
Since
\[
 \frac{2\tau^2}{a(a^2+\tau^2)}=\frac2a-\frac{2a}{a^2+\tau^2},
\]
the distribution kernel away from the diagonal is, up to the fixed Fourier
sign convention,
\begin{equation}
 K_\Gamma(u)=-\sum_{j\ge0}e^{-a_j|u|}
 =-\frac{e^{-|u|/2}}{1-e^{-2|u|}}.                            \tag{2}
\end{equation}
The omitted terms $2a_j^{-1}\delta_0$ are diagonal and therefore do not
enter an old/new cross block.

Let $I_N=(-T_N,T_N)$ and
$E_N=(T_N,T_{N+1})\cup(-T_{N+1},-T_N)$, with
$T_N=\tfrac12\log N$.

\begin{theorem}[Uniform Gamma old/new cross --- \PROVED]
The integral operator obtained by restricting (2) from $E_N$ to $I_N$ has
operator norm bounded by an absolute constant, uniformly in $N$.
\end{theorem}

\begin{proof}
For $0<u\le1$, (2) is bounded in modulus by
$C(1/u+1)$; for $u\ge1$ it is bounded by $Ce^{-u/2}$.  Near the right
endpoint write an old coordinate as $T_N-x$, $x>0$, and a newborn
coordinate as $T_N+y$, $0<y<T_{N+1}-T_N$.  The singular part is dominated
by $C/(x+y)$.  Extending both functions by zero to $\mathbb R_+$ reduces
the bound to the Carleman operator
\[
 Cf(x)=\int_0^\infty\frac{f(y)}{x+y}\,dy,
 \qquad \|C\|_{2\to2}=\pi.
\]
The constant and exponentially decaying remainders are bounded by Schur's
test.  The left endpoint is identical by reflection, while opposite
endpoints have distance $2T_N$ and give a smaller exponential bound.
Summing the four blocks proves a bound independent of $N$.
\end{proof}

\section{The two Tate channels}

The old/new blocks of a Tate rank-one channel are products of the
$L^2$ norms of $e^{t/2}$ or $e^{-t/2}$ on $I_N$ and on one component of
$E_N$.

\begin{theorem}[Uniform Tate old/new cross --- \PROVED]
Both primitive Tate channels have old/new cross norm bounded by an absolute
constant, uniformly in $N$.
\end{theorem}

\begin{proof}
For example,
\begin{align*}
 \|e^{t/2}\|_{L^2(I_N)}^2&=2\sinh T_N,\\
 \|e^{t/2}\|_{L^2(T_N,T_{N+1})}^2
 &=e^{T_N}(e^{\delta_N}-1),
 \qquad \delta_N=\frac12\log(1+N^{-1}).
\end{align*}
Their product is bounded because $e^{2T_N}=N$ and
$e^{\delta_N}-1=O(N^{-1})$.  On the opposite newborn component the
product is smaller.  Replacing $t$ by $-t$ interchanges the two endpoints,
so the same bound holds for the second Tate channel.
\end{proof}

Combining the two theorems, the complete affine Gamma--Tate old/new cross
$B_N^{\Gamma T}$ satisfies
\begin{equation}
                         \|B_N^{\Gamma T}\|\le C_{\Gamma T}   \tag{3}
\end{equation}
with an absolute constant.

\section{Cost in the normalized positive reference}

Row (d) proves for the positive reference
\begin{equation}
 \widehat{\mathcal R}_N\ge\alpha_N I,
 \qquad \alpha_N\sim\sqrt N.                                 \tag{4}
\end{equation}

\begin{corollary}[Gamma--Tate capacity is asymptotically negligible ---
\PROVED]
One has
\begin{equation}
 \bigl\|\widehat{\mathcal R}_N^{-1/2}B_N^{\Gamma T}\bigr\|^2
 \le\frac{C_{\Gamma T}^2}{\alpha_N}=O(N^{-1/2}).              \tag{5}
\end{equation}
If the normalized old defect obeys $D_{0,N}\ge\eta I$ for a fixed
$\eta>0$, then its additional Gamma--Tate return cost is
\begin{equation}
                         O(\eta^{-1}N^{-1/2}).                 \tag{6}
\end{equation}
\end{corollary}

\begin{proof}
Equation (5) follows immediately from (3)--(4).  Passing through the
normalized defect inverse costs at most $\|D_{0,N}^{-1}\|\le\eta^{-1}$,
which proves (6).
\end{proof}

The fixed-gap hypothesis in (6) must be propagated, not assumed.  The
strict arithmetic limit of v109--v111 suggests the induction
\begin{equation}
 \operatorname{Cap}_{N+1}
 \le s_*+o(1)+O(\eta^{-1}N^{-1/2}),
 \qquad
 s_*=\frac{\sqrt\pi}{2}e^{1/4}\operatorname{erf}\!\left(\frac12\right)
 <1.                                                          \tag{7}
\end{equation}
To make (7) a theorem about the signed physical Schur complement, one must
still verify that the Kato-transported residual innovation factors through
the physical delay isometry of v111 and the normalized defect operator,
with no additional noncontractive connection term.  This is precisely the
Gamma--Tate attachment (8)--(9) of v111.  Once that source identity is
proved, choose $\eta<(1-s_*)/3$ and then $N_0$ so large that both error
terms in (7) are below $\eta$; the fixed gap propagates for all
$N\ge N_0$.

\begin{center}
\begin{tabular}{>{\raggedright\arraybackslash}p{.55\linewidth}
                >{\centering\arraybackslash}p{.15\linewidth}
                >{\raggedright\arraybackslash}p{.20\linewidth}}
\toprule
Statement & Status & Location \\
\midrule
Uniform Gamma old/new cross & \PROVED & Theorem 1 \\
Uniform two-Tate old/new cross & \PROVED & Theorem 2 \\
Normalized Gamma--Tate cost $O(N^{-1/2})$ under a fixed defect gap & \PROVED & (5)--(6) \\
Physical defect-factor attachment & \OPEN & v111 (8)--(9) \\
Large-threshold gap induction & \OPEN & follows after that attachment \\
Condition $G$, row (d) & \OPEN & plus finite threshold certification \\
\bottomrule
\end{tabular}
\end{center}

\end{document}
